\chapter{When a potential is also a probability landscape}\label{ch:equilibrium-bridge}

The finite theory does not need a probability law. When a physical canonical model supplies one, however, the potential and relative surprisal become two expressions for the same landscape. We derive that bridge, including its units and measure, and distinguish it from the separate dynamics that may preserve a density.


\section{Two ways of learning about the same states}
A physical model can assign energy to a state. Repeated observations can estimate how often nearby states occur. In a properly specified thermal equilibrium benchmark, these two descriptions are related. This supplies a second mathematical reading of the same potential and, in applications, a way to compare physical and statistical descriptions.

The relationship is conditional. A useful inventory potential does not become a thermal energy by resemblance to a Gaussian curve. A stationary histogram does not establish detailed balance. We first state the density relation, then identify the additional assumptions required for its dynamical and thermodynamic interpretations.

\section{Canonical normalization}
Suppose a system has a specified energy $U(x)$, one fixed temperature $T>0$, an appropriate canonical ensemble, and a declared reference measure $\nu$. Assume its density relative to that measure is
\begin{equation}
 p(x)=Z_U^{-1}\exp[-\beta_{\mathrm{thermo}}U(x)],
 \qquad \beta_{\mathrm{thermo}}=\frac1{k_BT},
\end{equation}
with finite nonzero normalizer. Define the dimensionless potential
\begin{equation}\label{eq:thermal-potential}
 V(x)=\frac{U(x)-U(x^*)}{k_BT}.
\end{equation}
The constant reference energy can be absorbed into the normalizer, giving $p=Z^{-1}e^{-V}$.

Here $\beta_{\mathrm{thermo}}$ has units of inverse energy. It is not the dimensionless coefficient $\beta_{\mathrm{bridge}}$ used to compare an independently defined EBU potential with relative surprisal. Keeping those symbols separate prevents a normalization convention from looking like a fitted universal constant.

\begin{corollary}{The normalized bridge in its canonical scope}\label{cor:beta-one}
Under the canonical assumptions and Equation~\ref{eq:thermal-potential}, define $J_{\mathrm{prob}}(x)=-\ln p(x)$ and relative surprisal
$J_{\mathrm{rel}}(x)=J_{\mathrm{prob}}(x)-J_{\mathrm{prob}}(x^*)$. Then
\begin{equation}\label{eq:relative-surprisal}
 J_{\mathrm{rel}}(x)=V(x),\qquad \beta_{\mathrm{bridge}}=1.
\end{equation}
\end{corollary}
\begin{proof}
Taking logarithms gives $J_{\mathrm{prob}}=V+\ln Z$. Subtracting its value at $x^*$ removes the constant and uses $V(x^*)=0$.
\end{proof}

The equality is exact in this canonical normalization. It identifies the dimensionless bridge coefficient as one because the energy scale $k_BT$ has already been included in $V$. Comparing independently defined fields uses a different question: whether they admit a common physical normalization. Chapter~\ref{ch:moving-field} retains that distinction.

\section{Why an independent physical ruler matters}
An internally consistent score could still be multiplied by an arbitrary positive constant. Paths would still integrate, histories would still telescope and the signs of all coalition coefficients would remain unchanged. Those laws alone therefore cannot choose the size of the unit. Calibration enters because a denomination needs more than internal consistency.

The canonical benchmark supplies an independent physical scale: energy measured relative to the thermal energy $k_BT$. Under its assumptions, a difference of one in $V$ corresponds to a factor $e$ in the appropriate endpoint density ratio. This gives the number a meaning beyond an arbitrary label attached to a landscape.

For an independently proposed EBU field, the ambitious question is whether its physical construction agrees with the independently observed probability landscape at the predicted scale. Defining the potential from that same histogram and then recovering the histogram would only check an identity. Agreement between independent descriptions is the physically informative comparison.

This explains the role of the later E1a benchmark. A controlled optical system offers a clean setting in which field changes and normalization can be examined before heterogeneous environmental or economic actions are compared. Its importance lies in the denomination question, not in any claim that optical traps model an economy. The benchmark has not been executed by this integration.

A failure of the proposed denomination in its declared benchmark would challenge the broader claim. Success would support that limited bridge and motivate further work; it would not on its own establish a universal resource conversion. The theory here states exactly what the comparison means. Apparatus, test design and results belong to the later volumes.

\section{A density is relative to a measure}
For a continuous variable, the probability of an exact point is normally zero. A density tells us probability per unit of a declared measure. Saying a state has density twice another state's is meaningful only within that coordinate and measure convention.

If $y=f(x)$ is an invertible smooth change of coordinates and densities are reported relative to the new Lebesgue measure, then
\begin{equation}\label{eq:density-jacobian}
 p_y(f(x))=\frac{p_x(x)}{|\det Df(x)|}.
\end{equation}
A nonconstant Jacobian contributes a nonconstant log term. It need not cancel in an endpoint ratio or an interaction contrast. Transforming the reference measure consistently with the state description preserves the original scalar relation; silently changing the measure does not.

A density-of-states factor creates the same kind of issue. If an observed macrostate density is proportional to $g(x)e^{-V(x)}$, its negative log is $V(x)-\ln g(x)$ plus a constant. One must account for $g$ rather than announce a failure or confirmation of $J=V$ from a comparison using incompatible definitions.

\section{Curvature agrees more generally than covariance does}
Where the density relation holds smoothly in the declared coordinates,
\begin{equation}\label{eq:hessian-equality}
 \nabla^2J_{\mathrm{prob}}(x)=\nabla^2V(x)=H(x).
\end{equation}
This follows by differentiation of the log relation. It does not require a Gaussian density, and $H$ may vary with state.

A stronger result applies when $V=\tfrac12(x-x^*)^{\mathsf T}H(x-x^*)$ on all of $\mathbb R^d$, with constant symmetric positive-definite $H$, Lebesgue reference measure and no truncation. Then the normalized density is a global Gaussian with
\begin{equation}\label{eq:gaussian-covariance}
 \operatorname{Cov}(x)=\Sigma=H^{-1}.
\end{equation}
Completing the quadratic, or transforming to unit-normal coordinates, proves this covariance statement. It is not a general consequence of matching Hessians at one point.

Truncate the Gaussian to nonnegative stocks, impose a finite-mass boundary, or change the reference measure, and the global covariance need not equal the ambient inverse Hessian. A local quadratic approximation can still be useful, but it must be labeled as an approximation. In particular, the two-cell closed stock world of the opening chapters is not an unconstrained full-support Gaussian distribution.

\section{Conservation changes which directions can fluctuate}
If physical states satisfy $Ax=b$, an interior accessible displacement lies in the null space of $A$. Choose an orthonormal basis matrix $Q$ for that tangent space and write $x=x_c+Qy$. The restricted potential has Hessian $Q^{\mathsf T}H Q$ in the linear coordinates $y$.

Only in a full untruncated Gaussian on that affine subspace, with its induced measure and positive-definite restricted Hessian, does
\begin{equation}
 \operatorname{Cov}(y)=(Q^{\mathsf T}HQ)^{-1},\qquad
 \operatorname{Cov}(x)=Q(Q^{\mathsf T}HQ)^{-1}Q^{\mathsf T}
\end{equation}
follow. The ambient covariance is singular along conserved directions. For curved manifolds, coordinate curvature and the chosen intrinsic measure need explicit treatment; one should not casually apply this affine formula.

This is a practical protection against interpreting a conserved total as a failed fluctuation measurement. The inaccessible direction is not supposed to have the variance of an unconstrained system.

\section{The same density can support different dynamics}
A probability density tells us how mass is distributed across states, not how it circulates among them. For a diffusion with constant positive diffusion coefficient $D$, a drift of the form
\begin{equation}
 b=-D\nabla V+c,\qquad \nabla\!\cdot(c p)=0,
\end{equation}
can preserve $p\propto e^{-V}$ under appropriate boundary conditions while carrying stationary current $c p$. A nonzero current generally violates the reversible equilibrium interpretation for ordinary even state variables.

Appendix~\ref{app:probability} gives a rotating Gaussian example and a driven jump ring. They show why density agreement, stationarity, detailed balance, local detailed balance and entropy accounting must be treated as separate claims. A histogram alone cannot establish all of them.

\section{What this bridge offers EBU}
Within its canonical scope, the bridge gives EBU a physical and a probabilistic reading of the same state comparison. Energy reduction in units of $k_BT$ equals the logarithm of an endpoint density ratio. The unknown normalization cancels, and the curvature of relative surprisal equals the curvature of the potential in the stated coordinates.

The feedback tank illustrates why this is an additional branch of the theory. Its deterministic operating point, decaying modes and storage inequality specify no reservoir temperature, noise law or invariant density. They therefore supply no canonical normalization by themselves. A physical stochastic model may add such a structure, but it must be part of that model.

This precision makes the connection more useful to EBU. A system can use finite accounting for deterministic delivery and use the canonical bridge where a thermal model actually supplies it, without forcing all services into the same physical interpretation. The next chapter carries the density relation through the exact coalition transform.
